\documentclass[11pt]{article}

\usepackage[margin=1in]{geometry}
\usepackage[T1]{fontenc}
\usepackage[utf8]{inputenc}
\usepackage{lmodern}
\usepackage{booktabs}
\usepackage{tabularx}
\usepackage{array}
\usepackage{enumitem}
\usepackage{hyperref}

\hypersetup{
    colorlinks=true,
    linkcolor=black,
    urlcolor=blue,
    pdftitle={PANDA Project Overview},
    pdfauthor={AIMI 2526 Team}
}

\title{PANDA Project Overview\\
\large Proposal, Objectives, Code Overview, and Success Criteria}
\author{AIMI 2526 --- Radboud University}
\date{June 4, 2026}

\begin{document}
\maketitle

\section{Project Summary}

This project is a reproducible re-implementation of the PANDA (Prostate cANcer graDe Assessment) challenge pipeline for automated grading of prostate biopsy whole-slide images. The task is to predict the ISUP grade group from 0 to 5 directly from H\&E-stained pathology slides. The project is being carried out as the final project for the AI in Medical Imaging course at Radboud University.

The practical goal is not to reproduce the full complexity of the original winning Kaggle solution, but to build a clean, scientifically defensible system that captures the highest-impact ideas:

\begin{enumerate}[leftmargin=*]
    \item tile-based input instead of a single low-resolution thumbnail,
    \item a concat-tile-pooling model that aggregates features across tiles,
    \item fixed 5-fold cross-validation with out-of-fold (OOF) evaluation and ensembling.
\end{enumerate}

The public Kaggle leaderboard is no longer available for PANDA because the competition closed in 2020. For that reason, the project uses 5-fold cross-validation on the PANDA training set as its primary evaluation and external validation on the PBGG-1 dataset as its secondary evaluation.

\section{Problem Statement}

The core medical AI problem is the automated grading of prostate cancer severity from digitized biopsy slides. Each slide must be assigned an ISUP grade in $\{0,1,2,3,4,5\}$, where larger values correspond to more aggressive disease. The challenge is difficult because whole-slide images are extremely large, diagnostically relevant regions occupy only a fraction of the tissue, and the grade labels are ordinal rather than purely categorical.

The official PANDA metric is quadratic weighted Cohen's kappa (QWK), which rewards near-misses more than large-grade errors. This is appropriate for the clinical problem because predicting grade 3 instead of grade 4 is less severe than predicting grade 0 instead of grade 5.

\section{Project Proposal}

The proposal is to move from a simple baseline toward a more realistic whole-slide grading pipeline in staged steps.

\subsection{Baseline}

The starting point is a thumbnail-based model that takes a single resized $512 \times 512$ RGB image per slide and predicts ISUP grade with an EfficientNet backbone. This baseline is simple, cheap to run on Kaggle, and useful for validating the training and evaluation pipeline.

\subsection{Planned Improvement Path}

The intended improvement path is:

\begin{enumerate}[leftmargin=*]
    \item extract tissue-rich tiles from each whole-slide image,
    \item train a tile-aware model that pools tile features at the slide level,
    \item train the best model family across all 5 folds,
    \item ensemble full 5-fold model families fairly at OOF time,
    \item evaluate the final system on PBGG-1 for external validation.
\end{enumerate}

This design keeps the project manageable while still targeting the most important technical ideas from the PANDA literature and the winning competition solutions.

\section{Objectives}

\subsection{Primary Objective}

Build a reproducible pipeline that predicts ISUP grade from prostate biopsy whole-slide images and demonstrates strong performance under fair 5-fold validation.

\subsection{Secondary Objectives}

\begin{enumerate}[leftmargin=*]
    \item Compare multiple loss functions for thumbnail-based grading.
    \item Build and validate a tile extraction pipeline that can be reused across experiments.
    \item Train and evaluate a concat-tile-pooling model.
    \item Measure the benefit of fair OOF ensembling.
    \item Evaluate the final selected model family on PBGG-1.
    \item Produce a repo and notebook workflow that any teammate can rerun on Kaggle.
\end{enumerate}

\section{Success Criteria}

The project uses both performance and engineering success criteria.

\subsection{Performance Success}

\begin{itemize}[leftmargin=*]
    \item PANDA target: 5-fold mean QWK $\geq 0.83$.
    \item External validation target: PBGG-1 QWK $\geq 0.65$.
    \item Intermediate milestone: first tile-based fold-0 run should beat the clean thumbnail baseline and move toward $0.78+$ QWK.
\end{itemize}

\subsection{Engineering Success}

\begin{itemize}[leftmargin=*]
    \item The codebase supports baseline, tile-based, OOF, and ensemble evaluation from shared modules in \texttt{src/}.
    \item The tile dataset exists as a Kaggle-attachable artifact and is readable through \texttt{PandaTileDataset}.
    \item Each serious experiment is logged in \texttt{results.md}.
    \item The team can reproduce training and evaluation through committed notebooks and scripts.
\end{itemize}

\subsection{Scientific Success}

\begin{itemize}[leftmargin=*]
    \item Validation is fair and uses a fixed fold split.
    \item QWK is computed consistently from one shared implementation.
    \item Final reporting includes per-fold results, mean $\pm$ standard deviation, global OOF QWK, and external validation.
\end{itemize}

\section{Codebase Overview}

The repository is structured around reusable modules in \texttt{src/}, thin Kaggle notebooks, and a running experiment log.

\begin{table}[h]
\centering
\begin{tabularx}{\textwidth}{>{\ttfamily}l X}
\toprule
Path & Purpose \\
\midrule
src/eval.py & Shared evaluation utilities, including QWK and reporting helpers. \\
src/dataset.py & Thumbnail dataset and tile-aware dataset loading logic. \\
src/tiles.py & Tile extraction, tile scoring, and artifact serialization. \\
src/model.py & EfficientNet baseline and concat-tile-pooling model definitions. \\
src/train.py & Main training entry point for single-fold experiments. \\
src/oof.py & Fair out-of-fold evaluation for a full 5-fold model family. \\
src/inference.py & Model loading, prediction helpers, and ensemble prediction. \\
scripts/make_folds.py & One-time generation of the shared fold split. \\
scripts/preprocess_tiles.py & Full-dataset tile preprocessing and artifact writing. \\
notebooks/02\_train.ipynb & Kaggle wrapper for single-fold thumbnail training. \\
notebooks/02b\_train\_all\_folds.ipynb & Kaggle runner for training all folds. \\
notebooks/02c\_ensemble\_oof.ipynb & Fair OOF ensemble evaluation notebook. \\
notebooks/03a\_train0fold\_tiles.ipynb & Kaggle workflow for tiles and initial tile training. \\
notebooks/03\_external\_validation.ipynb & External validation workflow for PBGG-1. \\
results.md & Append-only experiment log. \\
\bottomrule
\end{tabularx}
\end{table}

\section{Technical Approach}

\subsection{Data Pipeline}

The baseline system uses a resized $512 \times 512$ thumbnail image per slide. The tile pipeline extends this by reading a whole-slide image or raster fallback image, padding it to the tile size, reshaping it into a grid, scoring tiles by tissue content, and keeping the top-$N$ tiles. The implementation supports saving tiles either as \texttt{.npy}/\texttt{.npz} arrays or vertically stacked \texttt{.png} files for inspection and Kaggle-friendly storage.

\subsection{Modeling}

Two model families are implemented:

\begin{itemize}[leftmargin=*]
    \item \textbf{Thumbnail baseline:} a single EfficientNet image encoder followed by a regression or ordinal head.
    \item \textbf{Concat-tile-pooling model:} tiles are encoded independently, reshaped to slide-level tile features, average pooled over tiles, and mapped to an ISUP prediction.
\end{itemize}

The current training code supports \texttt{efficientnet-b0} and \texttt{efficientnet-b1} backbones as well as three losses: SmoothL1, MSE, and ordinal binary cross-entropy on cumulative thresholds.

\subsection{Evaluation}

Evaluation is built around a fixed 5-fold split stored in \texttt{data/train\_folds.csv}. Full-family evaluation uses OOF prediction: each fold's validation slides are scored only by the model that was trained without that fold. This produces:

\begin{itemize}[leftmargin=*]
    \item per-fold validation QWK,
    \item mean $\pm$ standard deviation across folds,
    \item global OOF QWK pooled across all slides.
\end{itemize}

Ensembling is performed fairly by averaging predictions from full model families on the same validation fold. External validation is planned on PBGG-1 to assess generalization beyond PANDA.

\section{Current Project Status}

As of June 4, 2026, the codebase is substantially further than the original scaffold:

\begin{itemize}[leftmargin=*]
    \item Thumbnail training, OOF evaluation, and fair ensemble scaffolding are implemented.
    \item Tile extraction code and full-dataset preprocessing code are implemented.
    \item \texttt{PandaTileDataset} is implemented and can load \texttt{.npy}, \texttt{.npz}, or stacked \texttt{.png} tile artifacts.
    \item A first Kaggle tile dataset already exists for the team.
    \item The external validation notebook exists in the repository, but the final PBGG-1 evaluation workflow still needs to be completed and reported.
\end{itemize}

\section{Results So Far}

The main logged results to date are summarized in Table~\ref{tab:results}.

\begin{table}[h]
\centering
\caption{Current PANDA results from the experiment log.}
\label{tab:results}
\begin{tabularx}{\textwidth}{l l l l X}
\toprule
Setting & Input & Loss & Scope & QWK \\
\midrule
EfficientNet-B0 & thumbnail & SmoothL1 & fold 0 & 0.7117 \\
EfficientNet-B1 & thumbnail & SmoothL1 & fold 0 & 0.7125 \\
EfficientNet-B0 & thumbnail & MSE & fold 0 & 0.7123 \\
EfficientNet-B0 & thumbnail & ordinal BCE & fold 0 & 0.7314 \\
EfficientNet-B0 & thumbnail & SmoothL1 & clean 5-fold OOF & 0.7116 $\pm$ 0.0106, global OOF 0.7116 \\
EfficientNet-B0 & thumbnail & MSE & clean 5-fold OOF & 0.7030 $\pm$ 0.0055, global OOF 0.7030 \\
EfficientNet-B0 & thumbnail & ordinal BCE & clean 5-fold OOF & global OOF 0.7244 \\
\bottomrule
\end{tabularx}
\end{table}

The current best thumbnail-only family is therefore \texttt{efficientnet-b0 + ordinal}, with a global OOF QWK of 0.7244. This is a meaningful improvement over the SmoothL1 and MSE baselines, but it is still below the final project target and below the score range expected from successful tile-based PANDA pipelines.

The logged confusion behavior from the baseline experiments suggests that the model most often makes near-miss errors between adjacent grades and tends to under-call ISUP grade 5, which is a clinically important failure mode and one of the reasons tile-based localization is expected to help.

\section{Work Remaining}

\subsection{Track A: Data}

\begin{itemize}[leftmargin=*]
    \item Run the first full GPU tile-based fold-0 experiment.
    \item Log the first comparable tile QWK in \texttt{results.md}.
    \item Run the planned tile-count ablation at fixed tile size, for example $N \in \{16,32,36\}$.
    \item Confirm and document the final Kaggle dataset slug, format, and directory layout for the rest of the team.
    \item If raw TIFF access becomes practical, compare raw-slide tiles against the PNG fallback pipeline.
\end{itemize}

\subsection{Track B: Model}

\begin{itemize}[leftmargin=*]
    \item Treat \texttt{b0 + ordinal} as the current thumbnail reference.
    \item Run the first real tile model with \texttt{efficientnet-b0}.
    \item Compare tile configurations directly against the 0.7244 OOF thumbnail reference.
    \item If tiles clearly win, train the selected tile configuration across all 5 folds.
    \item Only after a stable tile baseline exists should optional techniques such as mixup be considered.
\end{itemize}

\subsection{Track C: Evaluation}

\begin{itemize}[leftmargin=*]
    \item Run a fair OOF ensemble using the best full-family checkpoints, especially SmoothL1 plus ordinal.
    \item Log family-level and ensemble-level OOF metrics in \texttt{results.md}.
    \item Complete the PBGG-1 evaluation workflow and report its QWK.
    \item Add test-time augmentation only after the final selected family or ensemble is fixed.
\end{itemize}

\section{Definition of Project Success}

This project will be considered successful if it satisfies the following four conditions:

\begin{enumerate}[leftmargin=*]
    \item \textbf{Code success:} the repository provides a clean end-to-end workflow for baseline, tiles, OOF evaluation, and final inference.
    \item \textbf{Experimental success:} the team demonstrates that its chosen tile-based model or ensemble improves on the thumbnail baseline under fair evaluation.
    \item \textbf{Performance success:} the final system approaches or exceeds the target of 0.83 mean QWK on PANDA and 0.65 QWK on PBGG-1.
    \item \textbf{Reporting success:} the final paper includes a coherent problem statement, method description, ablation table, and external validation discussion.
\end{enumerate}

\section{Conclusion}

The project already has a working baseline, fair OOF evaluation, tile extraction code, tile dataset loading, and an initial set of full-family thumbnail results. The most important remaining step is no longer building infrastructure from scratch; it is converting the existing tile pipeline into strong comparable GPU experiments and determining whether the tile-based model family can clearly surpass the current thumbnail ordinal reference. If that happens, the rest of the project becomes a focused evaluation and reporting problem rather than a pipeline-construction problem.

\end{document}
